\chapter{Four-Dimensional Supersymmetry and the Wess--Zumino Model}

\lectureoverview{The one-dimensional supersymmetry algebra of the preceding lecture involved only the Hamiltonian. Lorentz invariance now forces energy to rejoin the spatial momentum, and the supercharges must become two-component spinors. We reach that algebra by returning first to a massless Weyl fermion and its Hamiltonian \(H=\boldsymbol{\sigma}\cdot\mathbf p\). We then represent the supercharges as differential operators on four-dimensional superspace, isolate chiral superfields with the constraint \(\bar D_{\dot\alpha}\Phi=0\), and use superspace integration to construct a concrete field theory. The successive choices \(\Phi^\dagger\Phi\), \(\Phi^2\), and \(\Phi^3\) produce propagation, equal boson--fermion masses, Yukawa interactions, and a related scalar interaction. The auxiliary field \(F\) makes these relations transparent, while the closing loop diagrams show why supersymmetry softens ultraviolet sensitivity.}

\section{Why the Hamiltonian Cannot Stand Alone}

In the toy model, space was suppressed and the basic relation was
\begin{equation}
\{Q,Q^\dagger\}=2H.
\label{eq:l8-time-only-algebra}
\end{equation}
That form is adequate when time is the only coordinate. In a relativistic theory, however, \(H=P_0\) is one component of the four-momentum \(P_\mu\). A Lorentz boost mixes it with \(P_i\). An algebra whose right-hand side singled out \(P_0\) could not transform covariantly.

The repair must therefore do two things at once. It must replace \(H\) by all four components \(P_\mu\), and it must provide indices on the left that can be matched to a Lorentz vector on the right. The massless Dirac equation supplies the needed clue: a relativistic fermion may be described by a two-component spinor, and the matrices \(\sigma^\mu\) convert a vector index into a pair of spinor indices.

\section{Massless Fermions, Helicity, and Chirality}

Begin with the physical distinction between massive and massless particles. Take a massive spin-\(\tfrac12\) particle moving in the direction of its spin. Because its speed is less than the speed of light, an observer may overtake it. In the new frame the momentum reverses while the spin does not, so the helicity reverses. A Lorentz-invariant massive theory must consequently contain both helicities.

No observer can overtake a massless particle. Proper Lorentz transformations cannot reverse its momentum relative to its spin, so one irreducible massless representation may carry a single helicity. A spatial rotation can turn the flight direction around, but it does not turn a right hand into a left hand.

\begin{classroomqa}
\audiencequestion{Is the overtaking argument really about helicity rather than chirality? \lecturetimestamp{00:05:42}}
\lecturerresponse{Yes. Helicity is the projection of spin along momentum, whereas chirality is the eigenvalue of the chiral matrix, conventionally \(\gamma^5\). For a massless mode the two labels are tied together, which is why the language can sound interchangeable in this example. For a massive particle they are distinct.}
\end{classroomqa}

The state count requires one qualification.\editorialnote{A single massless Weyl field has one particle helicity and the opposite antiparticle helicity. Thus "half as many states" refers to the reduction from a four-component Dirac field to one two-component Weyl field; it does not mean that antiparticle states have disappeared.}

Write the Weyl field used in this opening argument as
\begin{equation}
\chi=
\begin{pmatrix}
\chi_1\\
\chi_2
\end{pmatrix},
\end{equation}
The lecture denotes this field by \(\psi_\alpha\); we call it \(\chi\) here so that it is not confused with the conjugate Weyl representation conventionally placed in a chiral superfield below.
\begin{equation}
H\chi=\boldsymbol{\sigma}\cdot\mathbf p\,\chi.
\label{eq:l8-weyl-hamiltonian}
\end{equation}
For momentum \(\mathbf p\), the matrix \(\boldsymbol{\sigma}\cdot\mathbf p\) has eigenvalues \(\pm|\mathbf p|\). The positive-frequency branch supplies the particle mode; the negative-frequency branch is reinterpreted through the antiparticle construction rather than retained as a spectrum unbounded below.

\begin{figure}[H]
\centering
\includegraphics[width=0.70\textwidth]{lecture_08_figure_02.png}
\lectureframe{00:07:44.271}
\caption{The chiral two-component spinor and the Hamiltonian \(H=\boldsymbol{\sigma}\cdot\mathbf p\). The final momentum symbol is partly covered, but the spoken equation fixes it.}
\end{figure}

The handedness label also depends on convention.\editorialnote{The board associates a handedness directly with the positive eigenvalue of \(\boldsymbol{\sigma}\cdot\mathbf p\). Which chirality receives that sign depends on the definitions of \(\sigma^\mu\), \(\bar\sigma^\mu\), and the Weyl field. We retain the eigenvalue statement but do not attach a convention-independent left/right label to it.}

Replacing \(H\) by \(i\partial_t\) and \(\mathbf p\) by \(-i\nabla\) gives
\begin{equation}
i\partial_t\chi=-i\boldsymbol{\sigma}\cdot\nabla\chi.
\end{equation}
Introduce
\begin{equation}
\begin{aligned}
\sigma^\mu&=(\mathbf 1,\boldsymbol{\sigma}),\\
\bar\sigma^\mu&=(\mathbf 1,-\boldsymbol{\sigma}),\\
\eta_{\mu\nu}&=\operatorname{diag}(1,-1,-1,-1).
\end{aligned}
\label{eq:l8-sigma-convention}
\end{equation}
Then the equation just obtained is
\begin{equation}
i\sigma^\mu\partial_\mu\chi=0.
\label{eq:l8-weyl-equation}
\end{equation}
The conjugate Weyl representation uses \(\bar\sigma^\mu\). It is this conjugate representation, named \(\psi_\alpha\) below, whose free equation is \(i\bar\sigma^\mu\partial_\mu\psi=0\). Proving covariance requires the spinor representation of the Lorentz group, but the important fact here is structural: the four matrices \(\sigma^\mu\) join one spacetime index to two spinor indices.

\section{The Four-Dimensional Superalgebra}

The odd coordinates must now transform like Weyl spinors:
\begin{equation}
\theta^\alpha=(\theta^1,\theta^2),
\qquad
\bar\theta^{\dot\alpha}=(\bar\theta^{\dot1},\bar\theta^{\dot2}).
\end{equation}
The bar denotes complex conjugation, and the dot distinguishes the conjugate Lorentz representation. There are two complex, or four real, Grassmann components. A general even superfield
\begin{equation}
\mathcal S=\mathcal S(x^\mu,\theta^\alpha,\bar\theta^{\dot\alpha})
\end{equation}
therefore has a finite Grassmann expansion. Its highest monomial is
\begin{equation}
\theta^2\bar\theta^2
\equiv
(\theta^1\theta^2)(\bar\theta^{\dot1}\bar\theta^{\dot2}),
\end{equation}
the product of all four independent odd coordinates. The coefficient of this monomial is commonly called the \(D\)-component.

The supercharges also carry spinor indices. In four-dimensional \(N=1\) supersymmetry their algebra is
\begin{align}
\{Q_\alpha,\bar Q_{\dot\beta}\}
&=2\sigma^\mu_{\alpha\dot\beta}P_\mu,
\label{eq:l8-superalgebra-main}\\
\{Q_\alpha,Q_\beta\}&=0,
&
\{\bar Q_{\dot\alpha},\bar Q_{\dot\beta}\}&=0,
\label{eq:l8-superalgebra-nilpotent}\\
[P_\mu,Q_\alpha]&=0,
&
[P_\mu,\bar Q_{\dot\alpha}]&=0,
&
[P_\mu,P_\nu]&=0.
\label{eq:l8-superalgebra-commutators}
\end{align}
The factor of two is conventional. The matrix
\(\sigma^\mu_{\alpha\dot\beta}\) contracts the vector index on \(P_\mu\) and leaves precisely the two spinor indices appearing on the left.

\begin{classroomqa}
\audiencequestion{Are \(\alpha\) and \(\beta\) on the sigma matrix its matrix indices? \lecturetimestamp{00:22:01}}
\lecturerresponse{Yes. More precisely, one is an undotted Weyl index and the other is a dotted conjugate index; \(\mu\) is the spacetime-vector index. The sigma matrix is the bridge between those descriptions.}
\end{classroomqa}

The cleaned notation makes that bridge explicit.\editorialnote{The lecture deliberately suppresses dotted indices and debates the order \(\alpha\beta\) versus \(\beta\alpha\) at the board. Equations~\eqref{eq:l8-superalgebra-main} and \eqref{eq:l8-superalgebra-commutators} restore the standard dotted/undotted notation without changing the argument being made.}

The last line says that the supercharges are conserved and compatible with translations. Momentum components commute because their derivative representatives commute. In the differential convention used below,
\begin{equation}
P_\mu=i\partial_\mu,
\end{equation}
so Eq.~\eqref{eq:l8-superalgebra-main} may equally be written
\begin{equation}
\{Q_\alpha,\bar Q_{\dot\beta}\}
=2i\sigma^\mu_{\alpha\dot\beta}\partial_\mu.
\end{equation}

\section{Supercharges as Superspace Derivatives}

A momentum translates an ordinary coordinate; a supercharge translates an odd coordinate while also shifting spacetime. Fix left Grassmann derivatives and define
\begin{align}
Q_\alpha
&=\partial_\alpha
-i\sigma^\mu_{\alpha\dot\alpha}
\bar\theta^{\dot\alpha}\partial_\mu,
\label{eq:l8-Q}\\
\bar Q_{\dot\alpha}
&=-\bar\partial_{\dot\alpha}
+i\theta^\alpha\sigma^\mu_{\alpha\dot\alpha}\partial_\mu.
\label{eq:l8-Qbar}
\end{align}
Acting with \(\epsilon^\alpha Q_\alpha+
\bar\epsilon_{\dot\alpha}\bar Q^{\dot\alpha}\) gives the infinitesimal supersymmetry variation of a superfield. Its components are simply the coefficients of the independent Grassmann monomials, and the differential operators mix those coefficients.

\begin{classroomqa}
\audiencequestion{Which \(x\)-derivative is meant in the shorthand \(\bar\theta\sigma\partial\)? \lecturetimestamp{00:25:40}}
\lecturerresponse{All four. The repeated \(\mu\) in \(\bar\theta^{\dot\alpha}\sigma^\mu_{\alpha\dot\alpha}\partial_\mu\) is summed, so each component of \(\sigma^\mu\) is paired with the matching derivative \(\partial/\partial x^\mu\).}
\end{classroomqa}

\begin{classroomqa}
\audiencequestion{Should the conjugate charge contain \(\theta\), not another \(\bar\theta\)? \lecturetimestamp{00:26:16}}
\lecturerresponse{Yes. \(Q\) differentiates with respect to \(\theta\) and contains \(\bar\theta\); \(\bar Q\) differentiates with respect to \(\bar\theta\) and contains \(\theta\). They are conjugate operators.}
\end{classroomqa}

\begin{classroomqa}
\audiencequestion{Is \(\sigma^\mu\) symmetric? \lecturetimestamp{00:27:04}}
\lecturerresponse{The Pauli matrices are Hermitian. Three of the four matrices in \((\mathbf1,\boldsymbol{\sigma})\) are symmetric, while \(\sigma^2\) is antisymmetric under ordinary transpose. Hermiticity, together with dotted and undotted index placement, is the relevant conjugation property.}
\end{classroomqa}

\begin{classroomqa}
\audiencequestion{Does the sigma matrix in the conjugate equation carry the indices in the opposite order? \lecturetimestamp{00:27:13}}
\lecturerresponse{Conjugation exchanges dotted and undotted representations, so the safe notation is \(\sigma^\mu_{\alpha\dot\alpha}\) and its conjugate \(\bar\sigma^{\mu\,\dot\alpha\alpha}\), rather than treating the two labels as interchangeable matrix slots.}
\end{classroomqa}

The formulas in Eqs.~\eqref{eq:l8-Q} and \eqref{eq:l8-Qbar} differ in superficial signs from the shorthand used at the board. Those signs are linked to whether \(P_\mu\) is \(+i\partial_\mu\) or \(-i\partial_\mu\), whether derivatives act from the left or right, and how conjugate odd variables are ordered. What matters is that one convention obeys the entire algebra at once.

\section{A Constraint That Survives Supersymmetry}

A general superfield contains many components. We seek a smaller class without destroying the symmetry. The lecture first gives an ordinary analogy. For a vector \(\mathbf V\), the constraint
\begin{equation}
V_x=0
\end{equation}
is not rotationally invariant: after a rotation, \(V_y\) contributes to the new \(x\)-component. By contrast,
\begin{equation}
\mathbf V^2=1
\end{equation}
is invariant and removes one independent degree of freedom. A useful constraint must map constrained fields into constrained fields under the symmetry.

Suppose a linear odd operator \(\bar D_{\dot\alpha}\) defines
\begin{equation}
\bar D_{\dot\alpha}\Phi=0.
\label{eq:l8-chiral-constraint}
\end{equation}
After a supersymmetry variation \(\delta\Phi=Q\Phi\), the same condition will hold if \(\bar D\) anticommutes with the supercharges:
\begin{equation}
\bar D(Q\Phi)=-Q(\bar D\Phi)=0.
\end{equation}
This is the abstract consistency test that motivates the construction, rather than a definition appearing without explanation.

\begin{classroomqa}
\audiencequestion{Why may \(Q\) and \(\bar D\) be switched in the last expression? \lecturetimestamp{00:36:56}}
\lecturerresponse{They are odd operators chosen to anticommute. Switching their order therefore introduces a minus sign; once \(\bar D\) reaches \(\Phi\), the original constraint makes the result zero.}
\end{classroomqa}

The required covariant derivatives are obtained by reversing the spacetime-derivative signs relative to the supercharges:
\begin{align}
D_\alpha
&=\partial_\alpha
+i\sigma^\mu_{\alpha\dot\alpha}
\bar\theta^{\dot\alpha}\partial_\mu,
\label{eq:l8-D}\\
\bar D_{\dot\alpha}
&=-\bar\partial_{\dot\alpha}
-i\theta^\alpha\sigma^\mu_{\alpha\dot\alpha}\partial_\mu.
\label{eq:l8-Dbar}
\end{align}
They satisfy
\begin{align}
\{D_\alpha,\bar D_{\dot\beta}\}
&=-2i\sigma^\mu_{\alpha\dot\beta}\partial_\mu,\\
\{D_\alpha,Q_\beta\}
&=\{D_\alpha,\bar Q_{\dot\beta}\}=0,\\
\{\bar D_{\dot\alpha},Q_\beta\}
&=\{\bar D_{\dot\alpha},\bar Q_{\dot\beta}\}=0.
\end{align}

\begin{classroomqa}
\audiencequestion{Is the only change from \(\bar Q\) to \(\bar D\) the sign of the spacetime-derivative term? \lecturetimestamp{00:37:58}}
\lecturerresponse{Exactly. That sign reversal makes the cross terms cancel when \(D\) is anticommuted with either supercharge.}
\end{classroomqa}

We impose only Eq.~\eqref{eq:l8-chiral-constraint}. Imposing both \(D_\alpha\Phi=0\) and \(\bar D_{\dot\alpha}\Phi=0\) would, through their anticommutator, force \(\partial_\mu\Phi=0\), leaving at most a spacetime-independent remnant.
\editorialnote{The board says that imposing both constraints would force \(\Phi=0\). The algebra strictly forces spacetime independence; additional boundary or component conditions may then reduce the field further.}

\section{Solving the Chiral Constraint}

Define the shifted coordinate
\begin{equation}
y^\mu=x^\mu+i\theta\sigma^\mu\bar\theta.
\label{eq:l8-chiral-coordinate}
\end{equation}
With the left-derivative convention above, a function
\begin{equation}
\Phi=\Phi(y,\theta)
\label{eq:l8-chiral-solution}
\end{equation}
obeys \(\bar D_{\dot\alpha}\Phi=0\). The chain rule makes the cancellation explicit. Differentiating through \(y^\mu\) gives the same \(i\theta\sigma^\mu\partial_\mu\Phi\) term that appears with the opposite sign in \(\bar D_{\dot\alpha}\).

Thus a chiral superfield has no independent \(\bar\theta\)-dependence. It may still contain \(\bar\theta\), but only through the combination \(y^\mu\). Its finite expansion is
\begin{equation}
\Phi(y,\theta)
=\phi(y)+\sqrt2\,\theta\psi(y)+\theta^2F(y).
\label{eq:l8-chiral-expansion-y}
\end{equation}
The three component types are a complex scalar \(\phi\), a Weyl fermion \(\psi_\alpha\), and a complex scalar \(F\). The last is auxiliary rather than a second propagating boson.

\begin{classroomqa}
\audiencequestion{Has \(x\) been replaced by a new coordinate \(y\) orthogonal to it? \lecturetimestamp{00:49:08}}
\lecturerresponse{No. \(y^\mu\) is shorthand for the particular combination in Eq.~\eqref{eq:l8-chiral-coordinate}; it is not an additional spacetime direction. Working in \(y\) makes chirality simple, and we translate back to \(x\) when constructing the ordinary field theory.}
\end{classroomqa}

Expanding Eq.~\eqref{eq:l8-chiral-expansion-y} at fixed \(x\) gives
\begin{align}
\Phi(x,\theta,\bar\theta)
={}&\phi
+\sqrt2\,\theta\psi
+\theta^2F
+i\theta\sigma^\mu\bar\theta\,\partial_\mu\phi
\notag\\
&-\frac{i}{\sqrt2}\theta^2
\partial_\mu\psi\,\sigma^\mu\bar\theta
+\frac14\theta^2\bar\theta^2\Box\phi.
\label{eq:l8-chiral-expansion-x}
\end{align}
All component fields on the right are evaluated at \(x\). The conjugate antichiral field depends on
\(\bar y^\mu=x^\mu-i\theta\sigma^\mu\bar\theta\) and \(\bar\theta\). A chiral superfield cannot consistently be real, because its conjugate satisfies the opposite chirality constraint.

\section{Superspace Integrals}

A full-superspace action has the form
\begin{equation}
S_D=\int d^4x\,d^2\theta\,d^2\bar\theta\;K,
\label{eq:l8-D-term-action}
\end{equation}
where \(K\) is a real superfield. Berezin integration selects its
\(\theta^2\bar\theta^2\) coefficient. After the odd integrations, only an ordinary spacetime Lagrangian remains.

There is a second invariant construction. If \(W(\Phi)\) is chiral, then
\begin{equation}
S_F=\int d^4x\,d^2\theta\;W(\Phi)
+\int d^4x\,d^2\bar\theta\;\overline{W}(\Phi^\dagger)
\label{eq:l8-F-term-action}
\end{equation}
is supersymmetric. The first integral selects the \(\theta^2\) component and the second its conjugate.

\begin{classroomqa}
\audiencequestion{Why is it legitimate to integrate only over \(d^2\theta\) when the expression has no \(\bar\theta\)-dependence? \lecturetimestamp{01:16:30}}
\lecturerresponse{Chirality is preserved by supersymmetry, and the highest \(\theta\)-component of a chiral superfield changes only by a spacetime total derivative. Its spacetime integral is therefore invariant, just as the full \(D\)-component integral is invariant.}
\end{classroomqa}

\section{The Free Chiral Multiplet}

Take the apparently structureless expression
\begin{equation}
S_{\mathrm{free}}
=\int d^4x\,d^2\theta\,d^2\bar\theta\;
\Phi^\dagger\Phi.
\label{eq:l8-free-superaction}
\end{equation}
The two factors depend on oppositely shifted coordinates:
\begin{equation}
\Phi=\Phi(x+i\theta\sigma\bar\theta,\theta),
\qquad
\Phi^\dagger
=\Phi^\dagger(x-i\theta\sigma\bar\theta,\bar\theta).
\end{equation}
Inside the spacetime integral we may shift the dummy variable \(x\). One shift can remove the Grassmann displacement from one factor, but then it doubles the relative displacement in the other. Taylor expansion in that nilpotent displacement terminates after finitely many terms.

This is where the derivatives come from. The superspace expression
\(\Phi^\dagger\Phi\) contains no explicit spacetime derivative, yet the shifted arguments generate first and second derivatives of the component fields. Multiplying the expansions and retaining only the top component gives, in the convention of Eq.~\eqref{eq:l8-sigma-convention},
\begin{equation}
\mathcal L_{\mathrm{free}}
=\partial_\mu\phi^\ast\partial^\mu\phi
+i\bar\psi\bar\sigma^\mu\partial_\mu\psi
+F^\ast F.
\label{eq:l8-free-component-lagrangian}
\end{equation}
Up to integration by parts, the scalar piece may instead be written schematically as \(-\phi^\ast\Box\phi\), which is the form approached in the board calculation.

\begin{classroomqa}
\audiencequestion{Where does the second line in the Taylor expansion of the superfield come from? \lecturetimestamp{01:03:00}}
\lecturerresponse{It comes from the component expansion \(\Phi(y,\theta)=\phi(y)+\sqrt2\theta\psi(y)+\theta^2F(y)\). Each component is then Taylor-expanded because \(y\) differs from \(x\) by a nilpotent Grassmann bilinear.}
\end{classroomqa}

\begin{classroomqa}
\audiencequestion{Should that last component carry \(\theta^2\)? \lecturetimestamp{01:04:54}}
\lecturerresponse{Yes. The auxiliary component is \(F\theta^2\). Since it already contains both independent \(\theta\)'s, any further Taylor shift carrying another \(\theta\) vanishes.}
\end{classroomqa}

\begin{figure}[H]
\centering
\includegraphics[width=0.82\textwidth]{lecture_08_figure_03.png}
\lectureframe{01:08:57.900}
\caption{The shifted-coordinate expansion being collected into scalar, fermion, and auxiliary terms of the ordinary Lagrangian.}
\end{figure}

The first two terms in Eq.~\eqref{eq:l8-free-component-lagrangian} yield the massless Klein--Gordon and Weyl equations. The last contains no derivative:
\begin{equation}
\frac{\partial\mathcal L}{\partial F}=F^\ast=0,
\qquad
\frac{\partial\mathcal L}{\partial F^\ast}=F=0.
\label{eq:l8-free-F-equation}
\end{equation}
It introduces no initial data and no on-shell particle.

\begin{figure}[H]
\centering
\includegraphics[width=0.82\textwidth]{lecture_08_figure_04.png}
\lectureframe{01:11:32.360}
\caption{The no-derivative auxiliary term \(F^\ast F\) at the end of the component Lagrangian.}
\end{figure}

The propagator language says the same thing. The inverse quadratic operator for \(F\) is momentum independent. Its Fourier transform is local:
\begin{equation}
\langle F(x)F^\ast(y)\rangle
\propto \delta^{(4)}(x-y).
\label{eq:l8-F-contact}
\end{equation}
An \(F\)-line may organize a diagram, but its endpoints coincide. It is a contact contraction, not the trajectory of a physical particle.

\section{The Quadratic Superpotential and Equal Masses}

Propagation came from a \(D\)-term. Mass comes from the chiral integral. Choose the conventional normalization
\begin{equation}
W_2(\Phi)=\frac12m\Phi^2.
\end{equation}
The \(\theta^2\) component is
\begin{equation}
\left.W_2(\Phi)\right|_{\theta^2}
=m\phi F-\frac12m\psi\psi.
\end{equation}
Together with the conjugate term, the component Lagrangian becomes
\begin{align}
\mathcal L
={}&
\partial_\mu\phi^\ast\partial^\mu\phi
+i\bar\psi\bar\sigma^\mu\partial_\mu\psi
+F^\ast F
\notag\\
&+m\phi F+m^\ast\phi^\ast F^\ast
-\frac12m\psi\psi
-\frac12m^\ast\bar\psi\bar\psi.
\label{eq:l8-mass-offshell}
\end{align}
The fermion bilinear is a Majorana mass for the Weyl field. The \(F\)-equations are algebraic:
\begin{equation}
F^\ast=-m\phi,
\qquad
F=-m^\ast\phi^\ast.
\label{eq:l8-mass-F-equations}
\end{equation}
Substitution gives
\begin{equation}
\begin{aligned}
\mathcal L_{\mathrm{on\ shell}}={}&
\partial_\mu\phi^\ast\partial^\mu\phi
+i\bar\psi\bar\sigma^\mu\partial_\mu\psi\\
&-|m|^2|\phi|^2\\
&-\frac12\left(m\psi\psi+m^\ast\bar\psi\bar\psi\right).
\end{aligned}
\label{eq:l8-mass-onshell}
\end{equation}
Thus the scalar mass and fermion mass are both \(|m|\). In the board's diagrammatic language, a scalar turns locally into \(F\) through one \(m\phi F\) insertion and immediately turns back through its conjugate. Two insertions supply \(|m|^2|\phi|^2\), exactly the scalar mass term found by solving Eq.~\eqref{eq:l8-mass-F-equations}.

The normalization can now be stated precisely.\editorialnote{The lecture intentionally suppresses numerical coefficients while multiplying superfields. The factors \(1/2\), \(1/3\), and \(\sqrt2\) used here are the conventional Wess--Zumino normalization; they make the component masses and couplings exact while preserving every qualitative step of the board derivation.}

\section{The Cubic Superpotential and Interactions}

Now add
\begin{equation}
W(\Phi)=\frac12m\Phi^2+\frac13g\Phi^3.
\label{eq:l8-superpotential}
\end{equation}
For a general holomorphic \(W\), the \(F\)-component is
\begin{equation}
\left.W(\Phi)\right|_{\theta^2}
=F W'(\phi)-\frac12W''(\phi)\psi\psi.
\label{eq:l8-W-components}
\end{equation}
Here
\begin{equation}
W'(\phi)=m\phi+g\phi^2,
\qquad
W''(\phi)=m+2g\phi.
\end{equation}
Before eliminating \(F\), the cubic term contains
\begin{equation}
g\phi^2F-g\phi\psi\psi+\text{h.c.}
\label{eq:l8-cubic-offshell}
\end{equation}
The second term is a Yukawa vertex joining one scalar to two fermion lines. The first joins two scalars to an auxiliary line.

\begin{figure}[H]
\centering
\begin{tikzpicture}[scale=0.95,>=Latex]
  \coordinate (vone) at (0,0);
  \draw[->,thick] (-2,1.1) -- (vone);
  \draw[->,thick] (vone) -- (-2,-1.1);
  \draw[thick,dashed] (vone) -- (2,0);
  \node[left] at (-2,1.1) {\(\psi\)};
  \node[left] at (-2,-1.1) {\(\psi\)};
  \node[right] at (2,0) {\(\phi\)};
  \node[below=8pt of vone] {\(-g\phi\psi\psi\)};
\end{tikzpicture}
\caption{Pedagogical redraw of the Yukawa interaction: two fermionic legs and one scalar leg meet with coupling proportional to \(g\).}
\end{figure}

The complete auxiliary-field equation is
\begin{equation}
F^\ast=-W'(\phi)=-(m\phi+g\phi^2).
\label{eq:l8-general-F-equation}
\end{equation}
Eliminating it produces the scalar potential
\begin{equation}
V(\phi)=|W'(\phi)|^2
=|m\phi+g\phi^2|^2.
\label{eq:l8-scalar-potential}
\end{equation}
Besides the mass term, this contains cubic scalar terms proportional to \(m^\ast g\) and a quartic term
\begin{equation}
V(\phi)\supset |g|^2|\phi|^4.
\label{eq:l8-quartic}
\end{equation}
In the auxiliary-line picture, two \(\phi^2F\) vertices are connected by the local \(F\)-contraction. Because Eq.~\eqref{eq:l8-F-contact} forces their positions to coincide, the result is precisely a four-scalar contact vertex.

\begin{figure}[H]
\centering
\begin{tikzpicture}[scale=0.95]
  \coordinate (vfour) at (0,0);
  \foreach \angle in {35,145,215,325}
    \draw[thick,dashed] (vfour) -- (\angle:1.9);
  \fill (vfour) circle (1.5pt);
  \node[below=8pt of vfour] {\(-|g|^2|\phi|^4\)};
\end{tikzpicture}
\caption{After the auxiliary field is eliminated, two local \(F\)-vertices become one four-scalar contact interaction.}
\end{figure}

Supersymmetry has done more than place a boson and a fermion in the same theory. It has forced their masses to agree, tied the Yukawa coupling to \(g\), and fixed the quartic scalar coupling to \(|g|^2\). These relations could be tested through masses, scattering amplitudes, and decays; they are not optional numerical coincidences.

\section{Paired Loop Effects}

The relations become most useful when interactions correct particle masses. The scalar self-energy receives a bosonic loop from the quartic vertex and a fermionic loop from two Yukawa vertices. Each carries two powers of the same coupling. A closed fermion loop has the relative minus sign required by Fermi statistics.

\begin{figure}[H]
\centering
\begin{tikzpicture}[>=Latex,scale=0.9]
  \begin{scope}[xshift=-3.3cm]
    \draw[thick,dashed] (-2,0)--(0,0)--(2,0);
    \draw[thick,dashed] (0,0) .. controls (-0.9,1.7) and (0.9,1.7) .. (0,0);
    \fill (0,0) circle (1.4pt);
    \node at (0,-0.55) {bosonic loop};
  \end{scope}
  \begin{scope}[xshift=3.3cm]
    \draw[thick,dashed] (-2,0)--(-0.8,0);
    \draw[thick,dashed] (0.8,0)--(2,0);
    \draw[->,thick] (-0.8,0) arc[start angle=180,end angle=0,x radius=0.8,y radius=0.8];
    \draw[->,thick] (0.8,0) arc[start angle=0,end angle=-180,x radius=0.8,y radius=0.8];
    \fill (-0.8,0) circle (1.4pt);
    \fill (0.8,0) circle (1.4pt);
    \node at (0,-1.35) {fermionic loop};
  \end{scope}
\end{tikzpicture}
\caption{The lecture's paired scalar self-energy mechanisms. Supersymmetry fixes the relative couplings, while the closed fermion loop supplies the opposite sign.}
\end{figure}

At one loop the leading ultraviolet-sensitive pieces cancel between the complete bosonic and fermionic contributions. The same symmetry organizes higher orders through Ward identities and nonrenormalization properties. The statement concerns symmetry-related sums, not a literal one-to-one cancellation of every graph; wave-function renormalization, finite effects, and other symmetry-preserving corrections can remain. This is the practical power of building the theory in superspace: the equality of masses and the relations among couplings are enforced before any individual Feynman integral is evaluated.\editorialnote{The board states the cancellation more broadly; the main text supplies the standard qualification.}

\section{What Comes Next}

The model constructed here is the simplest interacting chiral theory, commonly called the Wess--Zumino model. It is an example, not the most general supersymmetric field theory, but the pattern persists: represent fields by supermultiplets, write invariant superspace integrals, and let the Grassmann expansion enforce component relations.

Exact supersymmetry is not visible in the observed particle spectrum, so the next problem is how the symmetry can be broken while retaining enough of its ultraviolet control to be useful. The other promised direction is grand unification: how
\begin{equation}
SU(3)\times SU(2)\times U(1)
\end{equation}
may fit into a larger group such as \(SU(5)\), and how supersymmetric particle content changes that unification story. Those are the questions to which the remaining lectures turn.
